\section{A prime-only criterion for ordinary RH}
\label{sec:sr-prime}
Write \(\Lambda\) for the von Mangoldt function and
\(\gamma=\lim_{m\to\infty}(\sum_{r=1}^m r^{-1}-\log m)\).
For \(x\geq1\) put
\[
 T(x)=\sum_{n\leq x}\Lambda(n)(x/n-1),\qquad
 E(x)=T(x)-x\log x+(1+\gamma)x.
\]
All sums over \(n\) in this part are sums over positive integers. For an
integer \(m\geq1\) define
\[
 R(m)=\sum_{n\leq m}\Lambda(n)\left(1-\frac{m\bmod n}{n}\right),\quad
 P(m)=\sum_{\substack{p\leq m\\p\ {\rm prime}}}
 \log p\left(1-\frac{m\bmod p}{p}\right).
\]
The notation \(m\bmod n\) denotes the integer in \([0,n-1]\) congruent to
\(m\) modulo \(n\).

\begin{srtheorem}[Fourth-power prime criterion]\label{thm:sr-prime}
Suppose that for every \(\epsilon>0\) there is a constant
\(C_\epsilon\geq0\), independent of the subsequent interval, such that
arbitrarily large \(Y\) satisfy
\begin{equation}\label{eq:sr-prime}
 P(j^4)\geq\gamma j^4-C_\epsilon(j^4)^{1/2+\epsilon}
 \quad\hbox{whenever }j\in\mathbb N,\quad Y\leq j^4\leq Y^{3/2}.
\end{equation}
Then every nontrivial zero of the Riemann zeta function has real part
\(1/2\).
\end{srtheorem}
The hypothesis in this theorem is an arithmetic input. The theorem does not
assert it unconditionally. This section and the next two sections prove
the entire implication without importing a theorem from the Lean archive.
The criterion follows Anton's \texttt{FourthSamplesRH} and
\texttt{PrimeFourthSamples}; the exclusion of positive real zeros below
allows a shorter ordinary-zeta oscillation proof.

\subsection{Elementary bounds and the exact factorial identity}
Put \(\psi(x)=\sum_{n\leq x}\Lambda(n)\) and
\(\vartheta(x)=\sum_{p\leq x}\log p\). We first record a bound used several
times below:
\begin{equation}\label{eq:sr-cheb}
 0\leq\vartheta(x)\leq\psi(x)\leq3x\qquad(x\geq0).
\end{equation}
Indeed, for any integer \(a\geq1\), Legendre's formula gives
\[
 v_p\binom{2a}{a}=\sum_{r\geq1}
 \bigl(\lfloor2a/p^r\rfloor-2\lfloor a/p^r\rfloor\bigr).
\]
Every summand is nonnegative, and any \(p^r\in(a,2a]\) contributes one.
Consequently
\(\psi(2a)-\psi(a)\leq\log\binom{2a}{a}\leq2a\log2\).
Summing with \(a=1,2,\ldots\) along successive powers of two gives
\(\psi(2^k)\leq(2^{k+1}-2)\log2\). If \(x>1\), choose the least \(k\)
with \(x\leq2^k\); then \(2^k<2x\), and hence
\(\psi(x)<4x\log2<3x\). For \(0\leq x\leq1\) the sum is zero.
The elementary inequality \(4\log2<3\) follows, for example, by retaining
four positive terms in the series for \(\exp(3/4)>2\).

Since \(\log r=\sum_{n\mid r}\Lambda(n)\), interchanging two finite sums gives
\[
 \log(m!)=\sum_{n\leq m}\Lambda(n)\lfloor m/n\rfloor .
\]
Using \(m/n=\lfloor m/n\rfloor+(m\bmod n)/n\) now yields the exact identity
\begin{equation}\label{eq:sr-factorial}
 E(m)=\gamma m-R(m)+q(m),\qquad
 q(m)=\log(m!)-m\log m+m.
\end{equation}
Monotonicity of \(\log t\), applied on each unit interval, gives
\[
 m\log m-m+1\leq\log(m!)\leq m\log m-m+1+\log m.
\]
Thus \(1\leq q(m)\leq1+\log m\leq2\sqrt m\).
For the last inequality integrate \(1/t\leq1/\sqrt t\) from \(1\) to \(m\).
In particular,
\begin{equation}\label{eq:sr-integer}
 R(m)\geq\gamma m-E(m).
\end{equation}
A lower bound \(R(m)\geq\gamma m-Cm^{1/2+\epsilon}\) implies
\(E(m)\leq(C+2)m^{1/2+\epsilon}\).

Higher prime powers have a uniformly bounded cost:
\begin{equation}\label{eq:sr-discard}
 0\leq R(m)-P(m)\leq K\sqrt m,\qquad
 K=3+\frac{18}{\mathrm e\log2}.
\end{equation}
The first difference is nonnegative because every weight is in \([0,1]\).
For \(m\geq2\) it is at most
\[
 \psi(m)-\vartheta(m)
 =\sum_{r=2}^{\lfloor\log m/\log2\rfloor}\vartheta(m^{1/r})
 \leq3\sqrt m+\frac{3\log m}{\log2}m^{1/3}.
\]
The function \(v\exp(-v/6)\), for \(v\geq0\), has maximum \(6/\mathrm e\).
Taking \(v=\log m\) proves \eqref{eq:sr-discard}. The case \(m=1\) is zero.
It follows that the fourth-power criterion with \(R\) is equivalent to
\eqref{eq:sr-prime}, allowing a change of \(C_\epsilon\) by \(K\).

\subsection{Interpolation between fourth powers}
For any bounded interval of positive \(x\), the expression
\(T(x)=\sum_n\Lambda(n)(x/n-1)_+\) has only finitely many nonzero terms.
It is a convex continuous function. For \(0<a\leq x\leq b\), \(a<b\), put
\(\lambda=(b-x)/(b-a)\). Convexity and the identity defining \(E\) imply
\[
 E(x)\leq\lambda E(a)+(1-\lambda)E(b)
 +\lambda a\log a+(1-\lambda)b\log b-x\log x.
\]
The last three terms are at most
\((x-a)(b-x)/(2a)\). To see this, subtract that quadratic from the chord
of \(f(t)=t\log t\) minus \(f(x)\). The resulting function of \(x\) has
second derivative \(1/a-1/x\geq0\), and vanishes at both endpoints.
A convex function vanishing at both endpoints is nonpositive between them.
Therefore
\begin{equation}\label{eq:sr-interpolate}
 E(x)\leq\max(E(a),E(b))+\frac{(b-a)^2}{8a}.
\end{equation}
Take \(a=j^4,b=(j+1)^4\), with \(j\geq1\).
Then \(b\leq16a\) and
\(b-a=4j^3+6j^2+4j+1\leq15j^3\). The interpolation loss is at most
\(225\sqrt a/8\).

Fix \(\epsilon>0\) and an interval from \eqref{eq:sr-prime}.
Equation \eqref{eq:sr-factorial} gives
\(E(j^4)\leq(C_\epsilon+2)(j^4)^{1/2+\epsilon}\) at every supplied sample.
Set \(X=16Y\). For sufficiently large \(Y\),
\[
 16(16Y)^{5/4}\leq Y^{3/2}.
\]
If \(X\leq x\leq X^{5/4}\), its bracketing fourth powers satisfy
\(Y\leq a\leq x<b\leq Y^{3/2}\), since \(a\geq x/16\) and \(b\leq16x\).
Consequently \eqref{eq:sr-interpolate} gives throughout that interval
\[
 E(x)\leq
 \left((C_\epsilon+2)16^{1/2+\epsilon}+225/8\right)x^{1/2+\epsilon}.
\]
Increasing the lower threshold for \(Y\) absorbs the displayed constant
into \(x^\epsilon\). Thus for every \(\epsilon>0\) there are arbitrarily
late intervals
\begin{equation}\label{eq:sr-continuous}
 [X,X^{5/4}]\quad\hbox{on which}\quad E(x)\leq x^{1/2+2\epsilon}.
\end{equation}
Both the uniformity of \(C_\epsilon\) and the shortening of the interval
have been used explicitly.

\section{The zeta transform and its poles}\label{sec:sr-transform}
We give the analytic facts needed in the window argument. The underlying
tools are integration, Fourier series for a Gaussian, and Cauchy's
integral formula.

\subsection{Continuation, absence of positive real zeros, and reflection}
For \(\Re s>1\), the absolutely convergent series and Euler product give
\[
 \zeta(s)=\sum_{n\geq1}n^{-s}=\prod_p(1-p^{-s})^{-1},\qquad
 -\frac{\zeta'(s)}{\zeta(s)}=\sum_{n\geq1}\Lambda(n)n^{-s}.
\]
The product follows by expanding finitely many geometric series and then
using absolute convergence to pass to all primes. Its logarithm converges
absolutely, so the product is nonzero. Termwise differentiation of the
logarithm is locally uniform in this half-plane.

Partial summation gives
\begin{equation}\label{eq:sr-continuation}
 \zeta(s)=\frac{s}{s-1}
  -s\int_1^\infty \{t\}t^{-s-1}\,dt .
\end{equation}
The integral and all its derivatives converge locally uniformly for
\(\Re s>0\); its \(r\)-th derivative is dominated on
\(\Re s\geq a>0\) by a constant times
\(t^{-a-1}(1+\log t)^r\), which is integrable. This proves meromorphic
continuation there, with a simple pole of residue one at \(1\).
Furthermore
\[
 \int_1^m\{t\}t^{-2}\,dt=1+\log m-\sum_{r=1}^m1/r
\]
for integers \(m\geq1\). Taking the limit in this identity and then
expanding \eqref{eq:sr-continuation} at \(1\) shows
\begin{equation}\label{eq:sr-laurent}
 \zeta(s)=\frac1{s-1}+\gamma+O(s-1).
\end{equation}

The alternating series \(\eta(s)=\sum_{n\geq1}(-1)^{n-1}n^{-s}\) is
holomorphic for \(\Re s>0\). Summation by parts, using bounded partial
sums of the signs, bounds its tails uniformly on every compact subset
by a constant times \(M^{-a}\), where \(a\) is the minimum real part and
\(M\) the first tail index. In \(\Re s>1\), separating even and odd terms
gives \(\eta(s)=(1-2^{1-s})\zeta(s)\), and continuation gives the same
identity throughout \(\Re s>0\).
For real \(s>0\), pairing successive terms expresses \(\eta(s)\) as a
sum of positive terms, the first of which is \(1-2^{-s}>0\).
Thus \(\zeta(s)<0\) for \(0<s<1\); there are no positive real zeros.

There are no zeros on \(\Re s=1\) either. For \(\sigma>1\) and real \(t\),
the logarithmic Euler product and
\(3+4\cos u+\cos(2u)=2(1+\cos u)^2\geq0\) give
\[
 \zeta(\sigma)^3|\zeta(\sigma+it)|^4
 |\zeta(\sigma+2it)|\geq1.
\]
If \(t\ne0\) and \(\zeta(1+it)\) vanished to order \(m\geq1\), the left
side would be \(O((\sigma-1)^{4m-3})\) as \(\sigma\downarrow1\), since
\(\zeta\) is holomorphic at \(1+2it\). This tends to zero, a contradiction.
At \(t=0\) there is a pole. Hence every zero in \(\Re s>0\) lies in
\(0<\Re s<1\), is nonreal, and has finite positive multiplicity.
The zeros form a locally finite set: a nonzero holomorphic function has
isolated zeros, and a compact subset of its domain cannot contain an
infinite discrete zero set without an accumulation point.

For completeness, reflection can be recovered as follows.
The Gaussian Fourier transform is
\[
 \int_{\mathbb R}e^{-\pi t x^2}e^{-2\pi i\xi x}\,dx
 =t^{-1/2}e^{-\pi\xi^2/t}\qquad(t>0).
\]
The value at zero is the Gaussian integral, obtained by squaring and
integrating in polar coordinates. Differentiation in \(\xi\) and
integration by parts give the differential equation
\(F'(\xi)=-(2\pi\xi/t)F(\xi)\), proving the formula.
The smooth periodization \(\sum_{n\in\mathbb Z}e^{-\pi t(x+n)^2}\)
has these values at integral \(\xi\) as its Fourier coefficients.
Both the periodization and the resulting Fourier series converge
uniformly with all needed derivatives. Evaluation at \(x=0\) gives
\[
 \Theta(t)=\sum_{n\in\mathbb Z}e^{-\pi n^2t}
 =t^{-1/2}\Theta(1/t).
\]
Mellin integration, initially for \(\Re s>1\), and a split at \(t=1\)
now give
\begin{equation}\label{eq:sr-reflection}
 \pi^{-s/2}\Gamma(s/2)\zeta(s)
 =-\frac1s+\frac1{s-1}
 +\frac12\int_1^\infty(\Theta(t)-1)
 \left(t^{s/2-1}+t^{(1-s)/2-1}\right)\,dt .
\end{equation}
The integral is entire because of Gaussian decay, locally uniform with
every derivative. Its right side is invariant under \(s\mapsto1-s\).
This proves the functional equation and continuation to the whole plane.
The standard Euler product
\[
 \Gamma(z)^{-1}=z e^{\gamma z}
       \prod_{n=1}^{\infty}(1+z/n)e^{-z/n}
\]
shows that \(\Gamma\) has no zeros and has poles only at the nonpositive
integers. Its product converges locally uniformly by
\(\log(1+z/n)-z/n=O(n^{-2})\); the identity follows from Euler's
limit formula for the gamma integral. These elementary gamma facts and
\eqref{eq:sr-reflection} imply that a zero away from the negative even
integers and \(0\) reflects to a zero at \(1-s\).
At \(0\), comparison of the residues in \eqref{eq:sr-reflection} gives
\(\zeta(0)=-1/2\). Thus proving zero-freeness in \(\Re s>1/2\)
suffices for RH.

\subsection{The smoothed Mellin identity}
For \(\Re s>1\) and every positive integer \(n\), direct integration gives
\[
 \int_1^\infty(x/n-1)_+x^{-s-1}\,dx
 =\frac{n^{-s}}{s(s-1)}.
\]
Indeed the integral starts at \(n\), and integrating
\(n^{-1}x^{-s}-x^{-s-1}\) gives the displayed answer.
Absolute convergence justifies summing this identity after multiplication
by \(\Lambda(n)\): the sum of integrals of absolute values is
\(\sum_n\Lambda(n)n^{-\sigma}/(\sigma(\sigma-1))<\infty\),
where \(\sigma=\Re s>1\). Also
\(\int_1^\infty x\log x\,x^{-s-1}\,dx=(s-1)^{-2}\) and
\(\int_1^\infty x\,x^{-s-1}\,dx=(s-1)^{-1}\). Consequently
\begin{equation}\label{eq:sr-mellin}
 \Phi(s):=\int_1^\infty E(x)x^{-s-1}\,dx
 =-\frac{\zeta'(s)}{s(s-1)\zeta(s)}
       -\frac1{(s-1)^2}+\frac{1+\gamma}{s-1}.
\end{equation}
Initially this holds in \(\Re s>1\).
The expansion \eqref{eq:sr-laurent} cancels both apparent polar terms at
\(1\). Therefore \(\Phi\) is holomorphic on \(\Re s>0\) except at the
zeros of \(\zeta\). If \(\rho\) is a zero of multiplicity \(m\), factor
\(\zeta(s)=(s-\rho)^m h(s)\), \(h(\rho)\ne0\), locally. Then
\[
 \Phi(s)=\frac{c_\rho}{s-\rho}+H_\rho(s),\qquad
 c_\rho=-\frac{m}{\rho(\rho-1)}\ne0,
\]
where \(H_\rho\) is holomorphic near \(\rho\).
Every such pole is nonreal.

Set \(g(u)=E(e^u)\), \(u\geq0\). This is a continuous real function.
The trivial estimates \(\Lambda(n)\leq\log n\) and
\(\sum_{n\leq x}1/n\leq1+\log x\) give
\(T(x)\leq x\log x(1+\log x)\).
Thus \(|g(u)|\leq C_0 e^{2u}\) for a fixed \(C_0\):
each of \(1,u,u^2\), multiplied by \(e^{-u}\), is bounded on
\([0,\infty)\). The change of variables \(x=e^u\) gives
\[
 \Phi(s)=\int_0^\infty g(u)e^{-su}\,du\qquad(\Re s>2).
\]
For every \(k\geq0\), differentiation under the integral yields
\begin{equation}\label{eq:sr-deriv}
 A_k(s):=\frac{(-1)^k\Phi^{(k)}(s)}{k!}
 =\int_0^\infty g(u)\frac{u^k e^{-su}}{k!}\,du .
\end{equation}
On any closed sub-half-plane \(\Re s\geq2+a\), the differentiated
integrands are dominated by \(C_0u^k e^{-au}\), an integrable function.
This supplies the justification for every derivative in \eqref{eq:sr-deriv}.

\section{The window oscillation argument in full}\label{sec:sr-landau}
\begin{srtheorem}[One-sided windows exclude right-hand poles]\label{thm:sr-window}
Let \(0<\theta<1\) and \(\delta>0\). If arbitrarily large \(U\) satisfy
\[
 g(u)\leq e^{\theta u}\qquad(U\leq u\leq(1+\delta)U),
\]
then \(\zeta\) has no zero with real part greater than \(\theta\).
\end{srtheorem}

\subsection{A localized gamma kernel}
For \(a>0\), integration by parts \(k\) times gives
\[
 \int_0^\infty\frac{u^k e^{-au}}{k!}\,du=a^{-k-1},\qquad
 \frac{u^k e^{-au}}{k!}\leq a^{-k}.
\]
The second inequality follows by retaining the \(k\)-th term of the
series for \(e^{au}\). The boundary terms in the integration vanish
because exponential decay dominates each fixed power.

We will choose \(\sigma>2\), depending only on \(\delta\), and, for each
\(U\geq1\), an integer \(k\geq U\) such that
\begin{equation}\label{eq:sr-tail}
 \int_{(0,\infty)\setminus[U,(1+\delta)U]}
       \frac{u^k e^{(2-\sigma)u}}{k!}\,du
 \leq C_1\sigma^{-k}.
\end{equation}
Here is an explicit existence argument, including both tails.
Put \(\alpha=1+\delta/2\) and \(l=2+6/\delta\).
For sufficiently large \(\sigma>2+l\), both
\[
 -l(1+\delta)+\alpha\sigma
       \log\frac{\sigma}{\sigma-2-l}<-\tfrac12,\qquad
 l+\alpha\sigma\log\frac{\sigma}{\sigma-2+l}<-\tfrac12
\]
hold. The limits of their left sides as \(\sigma\to\infty\) are both
\(-1\), since \(\sigma\log(\sigma/(\sigma-a))\to a\).
Choose such a \(\sigma\), also with \(\alpha\sigma\geq1\), and set
\(k=\lceil\alpha\sigma U\rceil\).
For \(u\geq(1+\delta)U\), insert \(e^{l(u-(1+\delta)U)}\geq1\)
into the integrand. For \(0<u<U\), insert \(e^{l(U-u)}\geq1\).
Enlarging each resulting integral to \((0,\infty)\) gives the tail bound
\[
 \frac{e^{-l(1+\delta)U}}{(\sigma-2-l)^{k+1}}
 +\frac{e^{lU}}{(\sigma-2+l)^{k+1}}.
\]
Multiply by \(\sigma^k\).
In the first term, use \(k<\alpha\sigma U+1\), since its logarithmic
ratio is positive. In the second, use \(k\geq\alpha\sigma U\), since
its logarithmic ratio is negative. The sum is at most
\[
 e^{-U/2}\left(
 \frac{\sigma}{(\sigma-2-l)^2}+\frac1{\sigma-2+l}\right).
\]
This proves \eqref{eq:sr-tail} with the parenthesized expression as
\(C_1\). The same choice has \(k\geq U\).

\subsection{An inequality for the derivatives of the transform}
For a good interval \(I=[U,(1+\delta)U]\), let
\(w(u)=u^k e^{-\sigma u}/k!\).
For any real \(v\), \eqref{eq:sr-deriv} and the triangle inequality give
\[
 |A_k(\sigma+iv)|+A_k(\sigma)
 \leq\int_0^\infty(|g(u)|+g(u))w(u)\,du .
\]
The quantity \(A_k(\sigma)\) is real. On \(I\),
\(|g|+g=2\max(g,0)\leq2e^{\theta u}\); outside \(I\), it is at most
\(2C_0e^{2u}\). The kernel maximum and \eqref{eq:sr-tail} therefore give
\begin{equation}\label{eq:sr-key}
 |A_k(\sigma+iv)|+A_k(\sigma)
 \leq\frac{2\delta U}{(\sigma-\theta)^k}
       +\frac{2C_0C_1}{\sigma^k}.
\end{equation}
For every \(\eta>0\), all sufficiently large \(U\) satisfy
\(\delta U\leq e^{\eta U}\). Since \(\sigma-\theta<\sigma\),
\begin{equation}\label{eq:sr-key2}
 |A_k(\sigma+iv)|+A_k(\sigma)
 \leq C_2 e^{\eta U}(\sigma-\theta)^{-k},
 \qquad C_2=2+2C_0C_1.
\end{equation}
For every \(\eta>0\), this holds for arbitrarily large good \(U\),
with a corresponding \(k\geq U\), simultaneously for all \(v\).

\subsection{Isolating one pole}
Suppose for a contradiction that \(\rho_0\) is a zero with
\(\Re\rho_0>\theta\). Write \(Z\) for the zeros in the positive half-plane.
Consider the closed disk centered at the real point \(\sigma\), of radius
\(\sigma-\theta\). It lies in \(\Re s\geq\theta>0\) and contains finitely
many elements of \(Z\).

There exist a number \(D\leq\sigma-\theta\) and a zero \(\rho_1\) such that
the open disk \(B(\sigma,D)\) is pole-free and
\begin{equation}\label{eq:sr-horizontal}
 \sigma-\Re\rho_1<D.
\end{equation}
If the closed disk has no zeros, take \(D=\sigma-\theta,\rho_1=\rho_0\).
Otherwise choose a zero \(\rho_*\) in it closest to \(\sigma\) and take
\(D=|\sigma-\rho_*|,\rho_1=\rho_*\).
All positive-half-plane zeros are nonreal, so
\(\sigma-\Re\rho_*<|\sigma-\rho_*|\).
This proves \eqref{eq:sr-horizontal} in the second case as well.

We next find a center \(z=\sigma+iv\) and a radius \(r\leq D\) whose open
disk contains exactly one zero \(\rho\), with \(d=|z-\rho|<r\).
Start with \(z_0=\sigma+i\Im\rho_1\), so
\(|z_0-\rho_1|=\sigma-\Re\rho_1<D\).
Choose \(h>0\) so small that
\[
 |z_0-\rho_1|+2h<D,\qquad D+2h<\sigma .
\]
The set \(F=Z\cap\overline B(z_0,D+h)\) is finite. As \(v\) varies in
\((\Im\rho_1-h,\Im\rho_1+h)\), two distinct points \(a,b\) of \(F\)
can have equal distances from \(\sigma+iv\) for at most one value of \(v\).
Indeed equality of their squared distances is affine in \(v\).
If their imaginary parts coincide, equality would force equal real
parts too, since both real parts are less than \(1<\sigma\).
Exclude the finitely many possible tie values and choose \(v\) in the
remaining interval. There is a unique nearest point \(\rho\) of \(F\);
its distance is at most \(|z-\rho_1|<D\).
Every zero outside \(F\) has distance greater than \(D\) from \(z\),
by the triangle inequality.
Choose \(r\) strictly between \(d\) and the smaller of \(D\) and the
distances to all other points of \(F\). This proves the claim.
Both disks under discussion are contained in \(\Re s>0\).

\subsection{Cauchy estimates and the contradiction}
Choose \(b_0\) with \(d<b_0<r\).
Since \(\Phi\) is holomorphic in \(B(\sigma,D)\), Cauchy's formula
on the circle of radius \(b_0<D\) gives
\[
 |A_k(\sigma)|\leq M_0 b_0^{-k}
\]
with \(M_0\) independent of \(k\).
In \(B(z,r)\), subtract the single principal part
\(c_\rho/(s-\rho)\). The resulting function is holomorphic throughout
the disk, by the local factorization at \(\rho\). Cauchy's formula gives
\[
 \left|A_k(z)-\frac{c_\rho}{(z-\rho)^{k+1}}\right|
 \leq M_1 b_0^{-k}.
\]
Combining these estimates with \eqref{eq:sr-key2} yields, for arbitrarily
large good \(U\) and the corresponding \(k\geq U\),
\[
 \frac{|c_\rho|}{d^{k+1}}
 \leq(M_0+M_1)b_0^{-k}
       +C_2e^{\eta U}(\sigma-\theta)^{-k}.
\]
Here \(d>0\), since all zeros have real part less than \(1\), whereas
\(\Re z=\sigma>2\). Also \(d<b_0<r\leq D\leq\sigma-\theta\).
Multiply by \(d^k\), put \(\lambda=\log(b_0/d)>0\), and choose
\(0<\eta<\lambda\).
The right side is at most
\[
 (M_0+M_1)e^{-\lambda U}
   +C_2e^{-(\lambda-\eta)U},
\]
which tends to zero. The left side is the fixed positive number
\(|c_\rho|/d\). This contradiction proves
Theorem~\ref{thm:sr-window}.

\subsection{Completion of the conditional analytic deduction}
Suppose \(\zeta(\rho)=0\) with \(\Re\rho>1/2\).
The previous analytic facts give \(\Re\rho<1\).
Choose \(\epsilon>0\) so small that
\(\theta=1/2+2\epsilon<\Re\rho\).
Equation \eqref{eq:sr-continuous}, after \(x=e^u\), supplies arbitrarily
late intervals \([\log X,(5/4)\log X]\) with
\(g(u)\leq e^{\theta u}\).
Theorem~\ref{thm:sr-window}, with \(\delta=1/4\), rules out \(\rho\).
Thus there are no zeros in \(\Re s>1/2\).
Reflection from \eqref{eq:sr-reflection} proves
Theorem~\ref{thm:sr-prime}.
\qed
